\chapter{LITERATURE REVIEW}

This chapter provides a comprehensive review of prior research and existing engineering works in non-invasive mobile cardiovascular monitoring, digital signal processing, arrhythmia classification, dynamical cardiac rhythm modeling, and public consumer smartwatch validation. To address the fundamental objectives of our thesis, five directly related studies from the literature are critically analyzed, followed by a quantitative benchmark comparison against published public smartwatch PPG studies. For each core study, this review details: (1) the proposed methodology, (2) the datasets used, (3) the performance metrics evaluated, (4) the primary experimental findings, (5) the specific techniques that can be adopted (stolen) for our project, and (6) the fundamental differences and engineering advantages that distinguish our proposed mobile cardiac platform. Finally, a comparative synthesis table summarizes the state of the art relative to our work.

\section{Ballistocardiography and Seismocardiography: A Review of Recent Advances \cite{inan2015ballistocardiography}}

\textbf{Proposed Method:} Inan et al. \cite{inan2015ballistocardiography} presented a comprehensive engineering and physiological review of ballistocardiography (BCG) and seismocardiography (SCG), detailing how chest-mounted micro-electromechanical system (MEMS) accelerometers capture the micro-vibrations generated by myocardial contraction and blood ejection. The authors systematized the morphological landmarks of the dorso-ventral ($Z$-axis) SCG waveform---including mitral valve closure (MC), isovolumic moment (IM), aortic valve opening (AO), and aortic valve closure (AC)---and reviewed bandpass conditioning ($0.5\text{--}40\text{ Hz}$), adaptive filtering, and envelope-based systolic peak localization for heart rate and pulse rate variability (PRV) extraction.

\textbf{Data Used:} The review synthesized multi-center experimental cohorts spanning healthy resting adults, exercise recovery protocols, and heart-failure clinical studies recorded with lab-grade tri-axial MEMS accelerometers synchronized against simultaneous reference electrocardiography (ECG) and echocardiography.

\textbf{Performance Metrics Used:} Across the reviewed studies, key evaluation metrics included:
\begin{itemize}
    \item AO peak detection sensitivity ($\%$) and positive predictivity against ECG R-peaks.
    \item Heart rate estimation Mean Absolute Error ($\text{MAE}$ in bpm) and Bland--Altman $95\%$ limits of agreement.
    \item Pre-ejection period (PEP) and left ventricular ejection time (LVET) timing error ($\text{ms}$), plus Signal-to-Noise Ratio ($\text{SNR}$ in dB) across chest placement sites.
\end{itemize}

\textbf{Key Findings:} The authors established that placing a MEMS accelerometer along the dorso-ventral ($Z$-axis) over the lower sternum maximizes mechanical SNR and provides clear AO peak signatures corresponding closely to electrical QRS complexes. However, they identified three persistent barriers to consumer deployment: posture sensitivity (supine yields far higher waveform stability than standing), inter-subject anatomical variability, and severe vulnerability to gross body motion artifacts that overwhelm micro-$g$ cardiac accelerations.

\textbf{What We Can Steal (Adopt):} 
\begin{enumerate}
    \item \textit{Supine Lower-Sternum $Z$-Axis Placement Protocol:} We adopt the lower sternum (xiphoid region) supine $Z$-axis orientation as our primary sensing site to maximize AO peak SNR on consumer smartphones.
    \item \textit{Cardiac Mechanical Passband ($0.5\text{--}40\text{ Hz}$):} We adopt their recommended SCG frequency corridor via a zero-phase 4th-order Butterworth bandpass filter paired with a lightweight $200\text{ ms}$ moving-RMS energy envelope for unipolar AO peak extraction.
\end{enumerate}

\textbf{How We Are Different:} Most studies reviewed by Inan et al. relied on dedicated lab-grade wired accelerometers or custom chest patches taped to the sternum, paired with offline MATLAB post-processing. In contrast, our project uses an unmodified consumer smartphone's internal MEMS accelerometer (\$0 extra hardware), executes real-time zero-phase DSP with physiological energy gating and guided usability prompts in $< 500\text{ ms}$, and extends SCG monitoring to Two-Stage Hybrid Screening across the Rhythm-Disorder Cluster.

\newpage
\section{Cardiac Conduction Model for Generating 12-Lead Signals With Realistic Heart Rate Dynamics \cite{quiroz2018cardiac}}

\textbf{Proposed Method:} Quiroz-Ju{\'a}rez et al. \cite{quiroz2018cardiac} developed a biophysically grounded, heterogeneous nonlinear oscillator framework designed to simulate realistic 12-lead electrocardiograms incorporating natural heart rate variability (HRV), Mayer waves, and respiratory sinus arrhythmia (RSA). The authors modeled the natural cardiac pacemakers---the sinoatrial (SA) node, atrioventricular (AV) node, and His-Purkinje (HP) system---using modified Van der Pol oscillators coupled with time-delayed velocity terms. The depolarization and repolarization of atrial and ventricular myocardial tissues were simulated using coupled FitzHugh--Nagumo equations, building upon the foundational relaxation oscillation principles of van der Pol and van der Mark \cite{vanderpol1928}. Realistic rhythm variability was injected by modulating the SA node's intrinsic frequency parameter $f_1$, and a supervised perceptron-like learning algorithm fitted the model output to real subject recordings.

\textbf{Data Used:} The framework was fitted and validated against real 12-lead clinical ECG recordings from 13 human subjects (9 male, 4 female, aged 26--69) in the PhysioNet database, alongside documented pathological tracings (e.g., Wellens syndrome Type 1 ischemia and anterior myocardial infarctions).

\textbf{Performance Metrics Used:} The model was quantitatively evaluated using:
\begin{itemize}
    \item Mean Absolute Error ($\text{MAE}$) between simulated and empirical lead waveforms across 12 leads.
    \item Convergence rate of the perceptron fitting algorithm across optimization epochs.
    \item Spectral fidelity: power spectral density match across VLF, LF ($0.04\text{--}0.15\text{ Hz}$), and HF ($0.15\text{--}0.40\text{ Hz}$) frequency bands.
\end{itemize}

\textbf{Key Findings:} The model demonstrated that coupling nonlinear Van der Pol limit-cycle pacemakers with FitzHugh--Nagumo excitable media can accurately reproduce Einthoven's triangle lead geometries, realistic QT intervals, and complex pathological morphologies and arrhythmic rhythm perturbations. The perceptron algorithm achieved low MAE convergence within a finite number of iterations ($< 50$ epochs).

\textbf{What We Can Steal (Adopt):}
\begin{enumerate}
    \item \textit{Coupled Van der Pol Arrhythmia Formulation:} We adopt their coupled Van der Pol ODE structure (solved via Runge--Kutta RK45) to generate reproducible synthetic arrhythmic rhythm sequences with controlled ectopic perturbations, irregular AFib-like intervals, and rate shifts.
    \item \textit{PRV-Consistent Parameterization:} We map physiological heart rate and variability parameters into the oscillator's intrinsic frequency $\omega_0$ and damping $\mu$ so generated synthetic waveforms preserve realistic cardiac time scales.
\end{enumerate}

\textbf{How We Are Different:} Quiroz-Ju{\'a}rez et al. created an offline generative simulation tool fitted to 12-lead clinical ECG hardware. Rather than claiming a predictive 12-lead digital twin, our project repurposes the coupled Van der Pol ODE specifically as a \textbf{controlled arrhythmia simulation tool} to generate synthetic arrhythmic SCG test signals, bridging the scarcity of real pathological smartphone SCG recordings during end-to-end classifier verification.

\newpage
\section{Soft Real-Time Smartphone Cardiac Signal Processing \cite{tsamis2017soft}}

\textbf{Proposed Method:} Tsamis et al. \cite{tsamis2017soft} investigated the computational feasibility, latency constraints, and power efficiency of executing real-time cardiac signal filtering and QRS detection entirely on mobile smartphone hardware. The authors implemented a mobile processing pipeline comprising a bandpass integer filter, derivative operator, squaring function, and a moving-window integrator (the Pan--Tompkins algorithm \cite{pan1985real}). They analyzed performance trade-offs across native C/C++ implementations versus higher-level Java/Android frameworks, assessing thread contention between sensor I/O interrupts and DSP numerical computation.

\textbf{Data Used:} Benchmarks were performed using 30-minute standard clinical ECG records from the PhysioNet MIT-BIH Arrhythmia Database, streamed over Bluetooth Low Energy (BLE) to simulate real-time sensor ingestion, alongside synthetic test signals with controlled noise levels.

\textbf{Performance Metrics Used:} The study rigorously evaluated:
\begin{itemize}
    \item Execution latency per sample buffer ($\text{ms}$).
    \item CPU utilization percentage ($\%$) across single-core and multi-core configurations.
    \item Memory footprint ($\text{MB}$) and garbage collection interrupt frequencies.
    \item Processing throughput ($\text{samples/second}$).
    \item Battery consumption rate ($\%\text{ per hour}$).
\end{itemize}

\textbf{Key Findings:} The authors showed that smartphone processors easily support real-time cardiac filtering with latency under $10\text{ ms}$ per 1-second buffer when circular ring buffers are utilized. However, naive buffer re-allocation caused frequent garbage collection pauses, which introduced timing jitter into RR-interval measurements. The study demonstrated that fixed-size pre-allocated circular buffers are mandatory to maintain soft real-time execution guarantees.

\textbf{What We Can Steal (Adopt):}
\begin{enumerate}
    \item \textit{Pre-allocated Ring Buffers:} We adopt their recommendation of using fixed-length circular window buffers ($30\text{--}60\text{ s}$) in our mobile data acquisition layer, eliminating dynamic heap allocations that introduce timestamp jitter.
    \item \textit{Latency and Footprint Budgeting:} We adopt their execution efficiency protocol, setting a strict $< 500\text{ ms}$ total processing latency budget and $< 2\text{ MB}$ edge model footprint from raw buffer completion to full PRV extraction and rhythm screening.
\end{enumerate}

\textbf{How We Are Different:} Tsamis et al. required an external peripheral ECG hardware sensor communicating over Bluetooth, creating physical hardware dependencies. In our project, all signal acquisition is entirely self-contained within the smartphone, capturing sternal SCG vibrations via internal MEMS accelerometers, eliminating external sensor hardware completely.

\newpage
\section{A Real-Time Approach for Heart Rate Monitoring Using a Hilbert Transform in Seismocardiograms \cite{tadi2016realtime}}

\textbf{Proposed Method:} Tadi et al. \cite{tadi2016realtime} proposed a real-time heart rate monitoring pipeline directly on smartphone sternal seismocardiograms (SCG) acquired from a consumer mobile phone's internal 3-axis MEMS accelerometer placed on the chest. After bandpass filtering the dorso-ventral ($Z$-axis) acceleration, the authors computed the Hilbert transform to construct an analytic signal envelope, followed by low-pass smoothing and adaptive peak thresholding to localize aortic opening (AO) beats and estimate beat-to-beat heart rate.

\textbf{Data Used:} The method was validated on resting supine human subjects using simultaneous smartphone MEMS SCG recordings ($100\text{--}200\text{ Hz}$) synchronized with reference clinical electrocardiography (ECG) ground truth.

\textbf{Performance Metrics Used:} The evaluation focused on:
\begin{itemize}
    \item Beat detection Sensitivity ($\text{Sens} = \frac{\text{TP}}{\text{TP} + \text{FN}}$) and Positive Predictivity ($+\text{P} = \frac{\text{TP}}{\text{TP} + \text{FP}}$).
    \item Beat-to-beat heart rate Mean Absolute Error ($\text{MAE}$) and Root Mean Squared Error ($\text{RMSE}$) against reference ECG.
    \item Bland--Altman $95\%$ limits of agreement for interval estimation.
\end{itemize}

\textbf{Key Findings:} Tadi et al. demonstrated that an unmodified smartphone resting flat on the sternum can extract beat-to-beat heart rate with high clinical agreement against reference ECG ($> 98\%$ beat detection accuracy in stationary supine subjects). However, computing the full-window Hilbert transform introduces FFT/edge-ringing overhead on constrained mobile buffers, and the pipeline focused solely on healthy heart-rate estimation without automated physiological energy gating, usability evaluation, or rhythm-disorder classification.

\textbf{What We Can Steal (Adopt):}
\begin{enumerate}
    \item \textit{Smartphone Sternal SCG Acquisition:} We adopt their zero-hardware supine sternal MEMS SCG acquisition setup ($100\text{--}200\text{ Hz}$, $Z$-axis) and adaptive peak-detection workflow.
    \item \textit{Physiological Outlier Clamping:} We pair peak detection with explicit physiological interval bounds ($333\text{ ms} \le \text{PPI} \le 1500\text{ ms}$, corresponding to $40\text{--}180\text{ bpm}$) to prevent artifact spikes from corrupting PRV features.
\end{enumerate}

\textbf{How We Are Different:} Whereas Tadi et al. relied on the Hilbert transform envelope and stopped at healthy heart-rate estimation, our Chapter~3 pipeline replaces the Hilbert transform with a lightweight \textbf{$200\text{ ms}$ moving-RMS energy envelope} (zero FFT overhead, $< 500\text{ ms}$ total edge latency), adds an active physiological energy gate with measurable guided usability ($\text{SUS} \ge 75$, setup $< 60\text{ s}$), and extends the pipeline to Two-Stage Hybrid Screening across the Rhythm-Disorder Cluster plus controlled Van der Pol synthetic SCG validation.

\newpage
\section{Monitoring of Cardiovascular Condition via Mobile Phones Using Noise Variance \cite{aliev2012monitoring}}

\textbf{Proposed Method:} Aliev et al. \cite{aliev2012monitoring} proposed a robust mathematical framework for mobile cardiovascular monitoring based on statistical noise variance estimation. Recognizing that mobile cardiac recordings are frequently contaminated by uncontrollable environmental and physiological noise, the authors argued that attempting to clean corrupted signals with static filters often introduces deceptive waveform distortions. Instead, they developed an algorithm that continuously computes the cross-correlation variance of high-frequency background noise components. When the estimated noise variance $\sigma_{\varepsilon}^2$ exceeds an acceptable threshold, the signal segment is flagged as unviable, preventing false diagnostic outputs.

\textbf{Data Used:} The methodology was tested using digitized electrocardiograms from coronary care patients exhibiting varying degrees of myocardial ischemia and cardiac arrhythmias, superimposed with calibrated white Gaussian noise and muscular tremor artifacts.

\textbf{Performance Metrics Used:} Evaluated metrics included:
\begin{itemize}
    \item Signal-to-Noise Ratio ($\text{SNR}$ in dB) estimation accuracy.
    \item Diagnostic reliability index ($R_D$) comparing filtered noisy signals against true baseline records.
    \item False positive rate ($\text{FPR}$) of arrhythmia alarms under varying noise levels.
\end{itemize}

\textbf{Key Findings:} The authors proved that calculating noise variance prior to feature extraction drastically reduces false clinical alarms. Systems that attempted to detect peaks without variance gating exhibited false positive alarm rates exceeding $30\%$ in high-noise environments, whereas variance-gated processing reduced false alarms to below $4\%$.

\textbf{What We Can Steal (Adopt):}
\begin{enumerate}
    \item \textit{Absolute Physiological Energy Gating:} We directly adopt their noise variance principle in our Tier 2 (Stage 1) DSP pipeline by implementing an absolute energy gate ($0.005\text{ m/s}^2 \le \sigma_{raw} \le 0.50\text{ m/s}^2$) that immediately rejects stationary surface vibrations or excessive subject movement.
    \item \textit{Deterministic Signal Quality Index (SQI):} We use quantitative SNR estimation ($\ge 12\text{ dB}$) and variance boundaries as the deterministic gating trigger that confirms clinical validity before passing SCG intervals to hybrid rhythm screening.
\end{enumerate}

\textbf{How We Are Different:} Aliev et al. used noise variance purely as a passive offline metric on external multi-lead ECG recordings. In contrast, our platform uses signal quality gating as an active edge controller directly on consumer smartphone hardware: when raw sternal SCG fails physiological variance bounds or exhibits low SNR, the system immediately rejects corrupted segments and provides real-time posture and stillness prompts to guarantee artifact-free feature extraction.

\newpage
\section{Public Smartwatch PPG Benchmarks Versus Smartphone Sternal SCG}

To objectively evaluate the engineering trade-offs of using an unmodified smartphone rather than a dedicated consumer smartwatch, Table~\ref{tab:smartwatch_benchmark} synthesizes published validation benchmarks from large-scale public smartwatch studies alongside our proposed Smartphone Sternal SCG platform.

\begin{table}[!ht]
\footnotesize
\sloppy
\centering
\caption{Quantitative Comparison of Published Public Smartwatch PPG Benchmarks Versus Proposed Smartphone Sternal SCG}
\label{tab:smartwatch_benchmark}
\begin{tabular}{|p{2.5cm}|p{2.0cm}|p{1.4cm}|p{2.3cm}|p{2.4cm}|p{2.4cm}|}
\hline
\textbf{Device / Study} & \textbf{Modality \& Sensing Site} & \textbf{Extra Cost} & \textbf{Resting HR Accuracy (vs. ECG)} & \textbf{Arrhythmia / AFib Screening Metric} & \textbf{Key Physical Trade-off} \\ \hline
\textbf{Apple Watch} \cite{bent2020investigating, nelson2019accuracy, perez2019large} & Wrist optical PPG (+ optional 1-lead ECG) & \$249--\$399+ & $\text{MAE} \approx 2.0\text{--}4.4\text{ bpm}$ (rest); $1.8\text{--}5.9\text{ bpm}$ (24h) & Apple Heart Study ($n=419{,}297$): $\text{PPV} = 84\%$ for AFib on ECG patch & Continuous 24/7 wear; wrist PPG affected by skin tone, perfusion \& PTT jitter \\ \hline
\textbf{Fitbit} \cite{bent2020investigating, nelson2019accuracy, lubitz2022detection} & Wrist optical PPG & \$150--\$299+ & $\text{MAE} \approx 2.9\text{--}5.9\text{ bpm}$ (rest); $3.5\text{--}9.3\text{ bpm}$ (24h) & Fitbit Heart Study ($n=455{,}699$): $\text{PPV} = 98.2\%$ during concurrent ECG patch & Continuous passive tracking; lower resting optical precision under motion/dark skin \\ \hline
\textbf{Garmin} \cite{bent2020investigating} & Wrist optical PPG & \$200--\$500+ & $\text{MAE} \approx 3.4\text{--}6.5\text{ bpm}$ (rest) & Primarily fitness HR / abnormal HR threshold alerts & Optimized for sports/endurance; higher resting HR MAE than chest sensors \\ \hline
\textbf{Smartphone Sternal SCG (Proposed)} & \textbf{Lower-sternum MEMS SCG ($Z$-axis, supine)} & \textbf{\$0} (existing phone) & \textbf{Target: $\text{MAE} < 3.0\text{ bpm}$, $\text{RMSE} < 4.5\text{ bpm}$} & \textbf{Target: $\text{Sens} \ge 88\%$, $\text{Spec} \ge 85\%$, $F_1 \ge 0.86$ (Rhythm Cluster)} & \textbf{Central mechanical recoil (zero optical/PTT bias); resting supine use only} \\ \hline
\end{tabular}
\end{table}

Two critical engineering insights emerge from Table~\ref{tab:smartwatch_benchmark}. First, at rest, consumer wrist-worn PPG smartwatches exhibit resting heart-rate MAE between $2.0$ and $6.5\text{ bpm}$ \cite{bent2020investigating, nelson2019accuracy} while requiring \$150--\$500+ in dedicated hardware and remaining susceptible to optical skin-pigmentation attenuation and peripheral vascular Pulse Transit Time (PTT) jitter. By measuring central ventricular ejection recoil directly at the lower sternum, smartphone sternal SCG targets resting $\text{MAE} < 3.0\text{ bpm}$ at \textbf{\$0 incremental hardware cost} with zero optical distortion. Second, large-scale wearable trials such as the Apple Heart Study \cite{perez2019large} and Fitbit Heart Study \cite{lubitz2022detection} demonstrate that smartwatches excel at passive 24/7 ambulatory surveillance during daily movement. Consequently, rather than positioning smartphone sternal SCG as a replacement for continuous fitness wearables, our thesis also identifies an optional \textbf{Complementary Dual-Mode Pathway}: a smartwatch can flag ambulatory irregular-pulse episodes in daily life, while smartphone sternal SCG provides an accessible, zero-cost central mechanical confirmation check at rest.

\newpage
\section{Comparative Synthesis of Related Works}

Table~\ref{tab:literature_comparison} provides a comprehensive comparative matrix synthesizing the five reviewed core works alongside our proposed senior thesis platform.

\begin{table}[!ht]
\footnotesize
\sloppy
\centering
\caption{Comprehensive Comparison Between Related Works and Proposed Platform}
\label{tab:literature_comparison}
\begin{tabular}{|p{2.3cm}|p{2.1cm}|p{2.1cm}|p{2.2cm}|p{2.1cm}|p{2.2cm}|}
\hline
\textbf{Study / Work} & \textbf{Sensors \& Modality} & \textbf{Noise Gating \& Filtering} & \textbf{Classification / Target} & \textbf{Dynamical Simulation} & \textbf{Key Metrics Evaluated} \\ \hline
\textbf{Inan et al. (2015) \cite{inan2015ballistocardiography}} & Chest MEMS SCG / BCG & Bandpass ($0.5\text{--}40\text{ Hz}$), manual posture control & AO/AC valve timing \& HR/PRV review & None & AO Sensitivity (\%), HR MAE, PEP/LVET (ms) \\ \hline
\textbf{Quiroz-Ju{\'a}rez et al. (2018) \cite{quiroz2018cardiac}} & 12-lead ECG & Additive Gaussian, baseline wander & Waveform morphological fitting & Coupled Van der Pol + FHN ODEs & Waveform MAE across 12 leads \\ \hline
\textbf{Tsamis et al. (2017) \cite{tsamis2017soft}} & Wearable ECG (BLE stream) & Integer bandpass, Pan--Tompkins & Peak detection \& HR extraction & None & Latency (ms), CPU (\%), Memory (MB) \\ \hline
\textbf{Tadi et al. (2016) \cite{tadi2016realtime}} & Smartphone sternal MEMS SCG & Bandpass + Hilbert envelope (no energy gate) & Beat-to-beat HR monitoring only & None & Sensitivity (\%), $+$P (\%), HR MAE/RMSE \\ \hline
\textbf{Aliev et al. (2012) \cite{aliev2012monitoring}} & Mobile phone ECG & Noise variance ($\sigma^2$) gating & Cardiac status reliability index & None & False positive alarm rate (\%), SNR (dB) \\ \hline
\textbf{Our Proposed Platform} & \textbf{Sternal MEMS SCG (\$0 extra hardware; vs. public smartwatch PPG)} & \textbf{Zero-phase Butterworth (0.5--40 Hz), 200 ms Moving-RMS, Energy Gate \& Guided UI} & \textbf{Two-Stage Hybrid Screening: Rhythm-Disorder Cluster (Ectopic, AFib, Rate)} & \textbf{Controlled Van der Pol SCG Arrhythmia Simulator (RK45)} & \textbf{Acc/Sens/Spec/F1, MAE ($<3$ bpm), Latency ($<500$ ms), SUS ($\ge 75$)} \\ \hline
\end{tabular}
\end{table}

As summarized in Table~\ref{tab:literature_comparison}, prior works have addressed individual aspects of cardiac monitoring---such as chest-mounted SCG physiology \cite{inan2015ballistocardiography}, offline dynamical conduction modeling \cite{quiroz2018cardiac}, mobile computational efficiency \cite{tsamis2017soft}, smartphone SCG heart-rate tracking \cite{tadi2016realtime}, statistical noise gating \cite{aliev2012monitoring}, large-scale smartwatch PPG screening \cite{perez2019large, lubitz2022detection}, or connected IoT telemetry classification frameworks \cite{pandian2026internet}. However, prior approaches either rely on external bio-potential/wearable hardware, use heavier Hilbert envelopes without automated energy gating, omit formal usability validation, or stop at healthy heart-rate estimation without patient-stratified rhythm screening.

Our proposed project integrates these strengths into a unified, zero-hardware pipeline: it captures sternal SCG directly via the smartphone's internal MEMS accelerometer \cite{inan2015ballistocardiography, tadi2016realtime}, replaces the Hilbert transform with a lightweight $200\text{ ms}$ moving-RMS envelope reinforced by physiological energy gating and guided usability prompts \cite{aliev2012monitoring, brooke1996sus}, extracts verified clinical PRV metrics \cite{taskforce1996heart}, executes Two-Stage Hybrid Screening across the Rhythm-Disorder Cluster under strict patient-stratified benchmarks \cite{dechazal2004automatic, pandian2026internet}, benchmarks quantitatively against published public smartwatch PPG data \cite{bent2020investigating, nelson2019accuracy}, and employs a controlled Van der Pol oscillator simulator \cite{vanderpol1928, quiroz2018cardiac} to generate synthetic arrhythmic SCG test signals when real patient SCG data is unavailable.